\documentclass[12pt,a4paper]{article}

% --- Paquetes Básicos ---
\usepackage[utf8]{inputenc}
\usepackage[spanish]{babel}
\usepackage[margin=2.5cm]{geometry}
\usepackage{amsmath, amssymb, amsfonts}
\usepackage{graphicx}
\usepackage{hyperref}
\usepackage{enumitem}
\usepackage{titlesec}
\usepackage{cite}

% --- Configuración de Estilo ---
\titleformat{\section}{\large\bfseries}{\thesection.}{1em}{}
\hypersetup{colorlinks=true, linkcolor=blue, citecolor=blue, urlcolor=blue}

\title{\textbf{Proyecto de Investigación: Modelos Probabilísticos en Esports}\\ \Large Primera Entrega}
\author{Julian Jimenez y equipo}
\date{\today}

\begin{document}

\maketitle

\section{Tema y delimitación inicial}

\textbf{Tema general:} Modelado probabilístico y análisis de datos en deportes electrónicos (Esports).

\textbf{Fenómeno específico:} Inferencia bayesiana para la estimación de la probabilidad de victoria previa al inicio del encuentro (\textit{pre-match}) en el entorno competitivo de \textit{Counter-Strike: Global Offensive} (CS:GO). A diferencia de los modelos deportivos tradicionales que se centran en la clasificación binaria, este proyecto busca modelar la naturaleza estocástica del juego y la incertidumbre epistémica de las predicciones.

\textbf{Unidad de análisis:} Series profesionales oficiales (Best of 1, Best of 3) registradas en bases de datos integradoras como HLTV.org y PandasScore. El análisis se descompone en dos niveles: (1) probabilidad de victoria de la serie completa y (2) probabilidad específica por mapa individual, considerando las asimetrías tácticas de cada terreno de juego.

\textbf{Delimitación temporal y contextual:} La investigación se acota al ecosistema competitivo moderno de CS:GO/CS2 durante el periodo 2015--2026. Se priorizan equipos de nivel profesional (Tier 1 y Tier 2) debido a la alta fidelidad y disponibilidad de datos estructurados, lo que permite una convergencia robusta de los modelos de programación probabilística. El contexto se limita a la fase de \textit{pre-match}, omitiendo eventos \textit{in-play} para centrarse en la capacidad predictiva estratégica.

\section{Pregunta de investigación}

¿Cómo puede diseñarse un modelo de regresión logística bayesiana jerárquica que, al integrar el diferencial de ranking global, el rendimiento histórico segmentado por mapa y el factor de momento competitivo, logre una estimación calibrada de la probabilidad de victoria y una cuantificación explícita de la incertidumbre en partidas de CS:GO?

\vspace{1em}
\noindent \textbf{Declaración de tipo:} Esta es una \textbf{pregunta predictiva y mixta}. Es predictiva en su objetivo de evaluación fuera de muestra (out-of-sample) y mixta en su intención de explicar el peso específico (\textit{importance weights}) de cada variable táctica en la probabilidad final, ofreciendo una transparencia que los modelos de caja negra frecuentista no poseen.

\section{Mini-protocolo de revisión}

Para fundamentar técnicamente este proyecto, se diseñó un protocolo de revisión sistemática bajo el marco PICOC (Population, Intervention, Comparison, Outcome, Context), cuyos resultados se sintetizan a continuación:

\begin{itemize}[noitemsep]
    \item \textbf{Criterios de sistematicidad:} Se realizó una búsqueda estructurada en bases de datos indexadas (Scopus, IEEE Xplore, ACM y ScienceDirect) utilizando la cadena de búsqueda consolidada: \texttt{('esports' OR 'Counter-Strike') AND ('win probability' OR 'match prediction') AND ('Bayesian' OR 'probabilistic')}.
    \item \textbf{Criterios de inclusión:} Artículos empíricos (2015--2026) que implementen modelos predictivos en esports o deportes con similares dinámicas de equipo, y que reporten métricas de error probabilístico (Brier Score, Log-loss).
    \item \textbf{Criterios de exclusión:} Estudios no revisados por pares, análisis puramente teóricos de teoría de juegos sin datos reales, e investigaciones centradas exclusivamente en el comportamiento tóxico o psicología del jugador sin aporte estadístico al resultado del juego.
    \item \textbf{Hallazgos finales:} Se seleccionaron 12 artículos clave para la matriz de síntesis, identificando una clara tendencia hacia la clasificación binaria y una subutilización de la inferencia bayesiana en el pre-match.
\end{itemize}

\section{Síntesis crítica de la literatura}

La revisión de la literatura evidencia un crecimiento exponencial en la analítica de esports, pasando de descriptores estadísticos básicos a modelos complejos de aprendizaje profundo. Sin embargo, persisten patrones que limitan la aplicabilidad estratégica real:

\subsection{El paradigma dominante: Clasificación sobre Probabilidad}
La mayoría de los estudios representativos, como los de \textit{Smith et al. (2020)}, se enfocan casi exclusivamente en la métrica de \textbf{Accuracy} (Precisión). Al tratar la victoria como un evento categórico determinista, estos enfoques fallan en entornos donde el riesgo y la incertidumbre son fundamentales. Un modelo frecuentista que predice una victoria con un 51\% de probabilidad se reporta como un ``éxito'' si el equipo gana, ignorando que el modelo estaba casi en la zona de azar puro.

\subsection{Sesgos de Estacionariedad y la ``Ilusión del Tag''}
Un patrón recurrente identificado en modelos tipo TrueSkill o Glicko-2 (\textit{Muller et al., 2019}) es la asunción de que la fuerza de una organización (ej. FaZe, NaVi) es estable. La literatura ignora frecuentemente la volatilidad de los \textit{rosters}, un error crítico en CS:GO donde la sinergia de los 5 integrantes es el motor del desempeño. 

\subsection{Métricas de Evaluación: Brier Score y Calibración}
Solo estudios recientes de alta calidad, como el de \textit{Wang \& Davis (2023)}, introducen el \textbf{Brier Score} y el \textbf{Log-Loss} como métricas maestras. La síntesis crítica revela que, aunque los modelos de Redes Neuronales Profundas (\textit{Chen et al., 2022}) logran un alto ajuste, sufren de \textit{overconfidence} (exceso de confianza), asignando probabilidades extremas a eventos inciertos, lo cual es corregido en la literatura bayesiana mediante el uso de distribuciones a priori informativas.

\section{Identificación del vacío}

Basándose en la síntesis previa, se identifica un vacío de investigación clasificado como \textbf{metodológico-metodológico y empírico}. 

\textbf{Patrón identificado:} Las herramientas actuales priorizan el ``quién ganará'' sobre el ``qué grado de certidumbre tenemos''.

\textbf{Supuesto subyacente:} Se asume que el rendimiento histórico en mapas específicos es una variable estática o de segundo orden frente al ranking general.

\textbf{Limitación y detalle del vacío:} No se ha documentado un modelo de \textbf{Inferencia Bayesiana Jerárquica} que integre de manera transparente: (1) el diferencial de habilidad individual, (2) la especialización táctica por mapa y (3) la incertidumbre de datos escasos en equipos nuevos o niveles de competición inferiores. 

El vacío radica en el uso de \textbf{Intervalos de Credibilidad (HDI)} para la toma de decisiones. Mientras que la literatura clásica ofrece una probabilidad puntual (ej. 60\%), nuestro enfoque llena el vacío al proporcionar una distribución (ej. ``estamos 94\% seguros de que la probabilidad de victoria está entre 52\% y 68\%''). Este nivel de detalle es inexistente en los modelos populares y resulta vital para analistas profesionales de esports.

\section{Declaración formal de contribución}

Este proyecto contribuirá mediante el diseño de un sistema de \textbf{Inferencia Probabilística de Momento y Contexto (IPMC)}. 

\textbf{¿Qué haremos?} Desarrollaremos una arquitectura de programación probabilística en Python (PyMC) que trascienda la clasificación binaria tradicional. El proceso permitirá que la información ``fluya'' entre niveles (equipos, mapas y torneos), mejorando la predicción incluso cuando hay pocos datos de un enfrentamiento específico.

\textbf{¿Con qué datos?} Se integrarán datasets de HLTV y PandasScore (2020-2026), capturados mediante scripts propios que garantizan la trazabilidad. Utilizaremos las \textit{opening odds} (cuotas de apertura) de las casas de apuestas como distribuciones a priori informadas, una técnica sugerida por la literatura más avanzada pero raramente implementada.

\textbf{¿Con qué método?} Se implementará una regresión logística bayesiana jerárquica. Esto permitirá que la información "fluya" entre niveles (equipos, mapas y torneos), mejorando la predicción incluso cuando hay pocos datos de un enfrentamiento específico.

\textbf{¿Para qué evaluar?} Para establecer un nuevo estándar de \textbf{calibración predictiva} en el campo. Evaluaremos no solo si el modelo "acierta", sino si sus probabilidades son honestas respecto a la realidad empírica del juego.

\section{Alcance y viabilidad}

El proyecto se ejecutará en 12 semanas bajo un flujo de trabajo reproducible:

\begin{itemize}
    \item \textbf{Dataset:} 5,000+ partidas profesionales (CS:GO y CS2).
    \item \textbf{Variables:} Ranking Relativo, Winrate en Mapa (de los últimos 6 meses), y Factor de Fatiga (partidos jugados en la última semana).
    \item \textbf{Métricas:} \textbf{Expected Calibration Error (ECE)}, Brier Score y Log-loss.
    \item \textbf{Viabilidad:} Contamos con infraestructura de computación local y acceso a APIs de datos abiertos. La tecnología (Python/PyMC) es madura y el equipo tiene la competencia técnica para la implementación de algoritmos MCMC.
\end{itemize}

\textbf{No haremos:} Análisis de datos biométricos, predicción de victorias en tiempo real ni modelos para otros juegos de género distinto a los FPS tácticos.

\end{document}
